% GENERATED_BY: oc_core_monograph_translation_draft_pipeline_v1.py
% AI_ASSISTED_TRANSLATION_DRAFT: TRUE
% THEOREM_REVIEW_STATUS: PENDING
% THEOREM_REVIEW_MODE: FORMAL_AI_GATE__CODEX_CLI_CHATGPT
% THEOREM_REVIEW_REVIEWER_ID:
% THEOREM_REVIEW_AUTHORITY_CLASS:
% THEOREM_REVIEWED_AT:
% THEOREM_REVIEW_DOSSIER_REF:
% SOURCE_LANGUAGE: EN
% TARGET_LANGUAGE: RU
% SOURCE_EN_REF: content/k_levels/k12.tex
% ================================================================
% ==== FILE: content/k_levels/k12.tex
% ================================================================

% ==============================
%  Ontology of Continua — Core
%  K-level Module: K12 (Limit Continua)
%  FULL MODULE — FINAL
% ==============================

\subsubsection{$K_{12}$ Обзор}
\label{sec:k12-overview}

$K_{12}$ — это \emph{предельный континуум}: структурная неподвижная точка
рекурсивной метаиерархии, разработанной в $K_{11}$.
Это первый уровень, на котором:

\begin{itemize}
    \item рекурсивные метаструктуры стабилизируются,
    \item ветвление (Аксиома~21) достигает замыкания,
    \item существует глобальная структурная неподвижная точка,
    \item все операторы $(\Psi,\Phi,\Lambda,U,\Chi)$ действуют как симметрии,
    \item вся лестница $K_0\dots K_{11}$ становится представимой внутри одного
          инвариантного континуума.
\end{itemize}

Следовательно, $K_{12}$ — это не просто «следующий уровень», а структурный
предел K-башни в Core v2.6.

% ================================================================
\subsubsection{Пространство состояний $\Omega(K_{12})$}

\[
\Omega(K_{12}) =
\{
\Omega^*,\;
A^*,\;
P^*,\;
\Theta^*,\;
J^*,\;
C^*,\;
\mu^{(12)}_t
\},
\]

где:
\begin{itemize}
    \item $\Omega^*$ — предельное пространство состояний, содержащее все
          допустимые вложения $K_0\!\to\!K_{11}$,
    \item $A^*$ — глобальное множество осей, индуцированное замыканием всех
          осей предыдущих уровней,
    \item $P^*$ — потенциалы, переживающие рекурсивную сходимость,
    \item $\Theta^*$ — устойчивые пороги, сохраняемые под рекурсией,
    \item $J^*$ — потоки, остающиеся инвариантными относительно всех
          структурных операторов,
    \item $C^*$ — фиксированное семейство циклов, определяющее время $K_{12}$,
    \item $\mu^{(12)}_t$ — мера на предельных конфигурациях.
\end{itemize}

Все компоненты удовлетворяют глобальной инвариантности:
\[
X^* = \mathcal{R}(X^*),\quad \forall X \in \{\Omega, A, P, \Theta, J, C\}.
\]

% ================================================================
\subsubsection{Граница $\partial\Omega(K_{12})$}

Граница состоит из:
\begin{itemize}
    \item конфигураций, в которых рекурсия не сходится,
    \item точек, где ветвление порождает несовместимые глобальные топологии,
    \item состояний, нарушающих предельную инвариантность относительно
          операторов,
    \item расходимости циклов неподвижной точки,
    \item недопустимости в $M_{12}$.
\end{itemize}

Пересечение $\partial\Omega(K_{12})$ разрушает предельную структуру и
возвращает систему к неудачному режиму $K_{11}$.

% ================================================================
\subsubsection{Оси $A(K_{12})$}

\[
A(K_{12}) =
\overline{\bigcup_{i=0}^{11} A(K_i)}
\quad \text{subject to recursive closure}.
\]

Это множество осей:
\begin{itemize}
    \item замкнуто относительно рекурсии из $K_{11}$,
    \item устойчиво относительно всех операторов $(\Psi,\Phi,\Lambda,U,\Chi)$,
    \item минимально среди множеств осей, удовлетворяющих глобальной
          инвариантности.
\end{itemize}

Каждая ось становится \emph{направлением симметрии} в $K_{12}$.

% ================================================================
\subsubsection{Потенциалы $P(K_{12})$}

\[
P(K_{12}) = \lim_{n\to\infty} P^{(n)}(K_{11}),
\]

то есть потенциалы, которые остаются конечными, согласованными и замкнутыми
относительно рекурсии:
\begin{itemize}
    \item $P_{\mathrm{fix}}^*$ — потенциал, поддерживающий неподвижную точку,
    \item $P_{\mathrm{coh}}^*$ — потенциал глобальной когерентности,
    \item $P_{\mathrm{branch}}^*$ — стабилизированный потенциал ветвления,
    \item $P_{\mathrm{embed}}^*$ — универсальный потенциал вложения.
\end{itemize}

% ================================================================
\subsubsection{Пороги $\Theta(K_{12})$}

\[
\Theta(K_{12}) =
\lim_{n\to\infty} \Theta^{(n)}(K_{11}),
\]

давая:
\begin{itemize}
    \item $\Theta^*_{\mathrm{fix}}$ — порог существования неподвижной точки,
    \item $\Theta^*_{\mathrm{coh}}$ — порог глобальной когерентности,
    \item $\Theta^*_{\mathrm{branch}}$ — порог замыкания ветвления,
    \item $\Theta^*_{\mathrm{inv}}$ — порог инвариантности операторов,
    \item $\Theta^*_{\mathrm{collapse}}$ — предел структурного коллапса.
\end{itemize}

Все пороги становятся \emph{глобально устойчивыми константами} для
предельного континуума.

% ================================================================
\subsubsection{Потоки $J(K_{12})$}

\[
J(K_{12}) = \{J^*\mid \Phi(J^*)=J^*\}.
\]

Следовательно, каждый поток должен быть:
\begin{itemize}
    \item инвариантным относительно рекурсивной эволюции,
    \item совместимым с предельными порогами,
    \item глобально когерентным на всех проекциях и вложениях,
    \item фиксированным под действием оператора.
\end{itemize}

Это единственный уровень, на котором \emph{все потоки являются потоками
симметрии}.

% ================================================================
\subsubsection{Циклы $C(K_{12})$}

\[
C(K_{12}) = \{ C^*\mid \Pi(C^*)>\Theta_{\mathrm{time}}^*,\;
\Lambda(C^*)=C^*\}.
\]

Это циклы, которые:
\begin{itemize}
    \item сохраняются через все рекурсивные уточнения,
    \item определяют возникающее время $K_{12}$,
    \item остаются инвариантными относительно ветвления и действия операторов,
    \item служат структурным каркасом континуума.
\end{itemize}

% ================================================================
\subsubsection{Time $\tau(K_{12})$}

Время определяется инвариантными циклами:
\[
\tau(K_{12}) = \max_j \tau_{C_j^*}.
\]

Свойства:
\begin{itemize}
    \item все временные направления являются симметриями,
    \item не могут возникать новые временные направления,
    \item время глобально когерентно и инвариантно относительно рекурсии.
\end{itemize}

% ================================================================
\subsubsection{Континуумность $k(K_{12})$}

\[
k_{12} =
H(\Omega(K_{12}))
\cdot
\frac{|C^*|}{|C_{\max}|}
\cdot
\frac{|A^*|}{|A_{\max}|}
\cdot
\frac{P_{\mathrm{fix}}^*}{P_{\mathrm{fix,max}}}
\cdot
\left(1-\frac{T_{12}}{\Theta^*_{\mathrm{collapse}}}\right)_+.
\]

$k_{12}$ достигает максимально возможного значения среди всех уровней $K$:
\[
k_{12} = \max_{i=0\dots12} k(K_i).
\]

% ================================================================
\subsubsection{Структурное напряжение $T(K_{12})$}

Источники:
\begin{itemize}
    \item неразрешенные несовместимости ветвления,
    \item сбои инвариантности операторов,
    \item нестабильность неподвижных точек (редкая, но возможная),
    \item несогласованности проекции–вложения,
    \item нарушения ограничений $M_{12}$.
\end{itemize}

Коллапс происходит, если:
\[
T(K_{12}) > \Theta^*_{\mathrm{collapse}}.
\]

% ================================================================
\subsubsection{Энергия $E(K_{12})$}

\[
E(K_{12}) =
E_{\mathrm{fix}}^*
+E_{\mathrm{coh}}^*
+E_{\mathrm{branch}}^*
+E_{\mathrm{embed}}^*
+E_{\mathrm{sym}}.
\]

Интерпретация:
\begin{itemize}
    \item минимальная энергия, требуемая для сохранения структуры неподвижной
          точки,
    \item энергия поддержания глобальной когерентности,
    \item энергия, стабилизирующая замыкание ветвления,
    \item энергия, поддерживающая универсальные вложения,
    \item энергия симметрии инвариантности операторов.
\end{itemize}

$K_{12}$ — это наиболее энергоустойчивый допустимый континуум.

% ================================================================
\subsubsection{Операторы на $K_{12}$ ($\Psi$, $\Phi$, $\Lambda$, $U$, $\Chi$)}

На предельном уровне:

\begin{itemize}
    \item $\Psi$ — становится тождеством на допустимых состояниях,
    \item $\Phi$ — становится глобальным потоком инвариантности,
    \item $\Lambda$ — замыкание всех предельных циклов,
    \item $U$ — универсальный оператор вложения,
    \item $\Chi$ — стабилизированный оператор ветвления.
\end{itemize}

Алгебра операторов становится абелевой с точностью до классов эквивалентности:
\[
[\Psi,\Phi]=[\Phi,\Lambda]=\dots = 0.
\]

% ================================================================
\subsubsection{Процессы на $K_{12}$}

\begin{itemize}
    \item сходимость рекурсивных структур,
    \item замыкание топологий ветвления,
    \item формирование глобальных неподвижных точек,
    \item коллапс недопустимых иерархий,
    \item возникновение полной симметрии под действием операторов,
    \item стабилизация времени и циклической структуры.
\end{itemize}

% ================================================================
\subsubsection{Предсказания для $K_{12}$}

\begin{enumerate}
    \item Любая допустимая рекурсивная метаиерархия в конечном итоге
          сходится к структурам, представимым в $K_{12}$.
    \item За пределами этого уровня не могут возникать новые оси или
          потенциалы.
    \item Все действия операторов сводятся к глобальным симметриям.
    \item Ветвление становится топологически замкнутым; новые ветви не
          возникают.
    \item Неспособность достичь глобальной неподвижной точки вынуждает
          коллапс обратно в $K_{11}$.
\end{enumerate}

% ================================================================
\subsubsection{Эксперименты для $K_{12}$}

Существуют только косвенные эксперименты:
\begin{itemize}
    \item симуляции рекурсивной сходимости,
    \item вычисление неподвижной точки в метапространствах высокого порядка,
    \item тесты инвариантности под алгеброй операторов,
    \item тесты устойчивости для замыканий ветвления,
    \item моделирование допустимости в рамках $M_{12}$.
\end{itemize}

% ================================================================
\subsubsection{Коллапс и смерть $K_{12}$}

Смерть, если:
\[
\Omega(K_{12})=\varnothing.
\]

Механизмы:
\begin{itemize}
    \item рекурсивная дивергенция на глубоких уровнях,
    \item нарушения инвариантности операторов,
    \item неудача замыкания ветвления,
    \item глобальная некогерентность неподвижных точек,
    \item несовместимость с метапространством $M_{12}$.
\end{itemize}

% ================================================================
\subsubsection{Фальсифицируемость $K_{12}$}

Опровержимо, если:
\begin{itemize}
    \item метаиерархия сходится без формирования инвариантных неподвижных точек,
    \item ветвление остается незамкнутым, но устойчивость сохраняется,
    \item после замыкания возникают новые оси или потенциалы,
    \item алгебра операторов не становится симметричной,
    \item допустимая рекурсия порождает несходящиеся иерархии.
\end{itemize}

% ================================================================
\subsubsection{Ветвление / онтологическое положение $K_{12}$}

Структура ветвления:
\begin{itemize}
    \item все ветви замыкаются,
    \item новые ветви не формируются,
    \item топология становится глобально компактной и инвариантной.
\end{itemize}

Онтологическая позиция:
\[
K_{11} \;\longrightarrow\; K_{12} \quad \text{(предельный континуум)}.
\]

За пределами $K_{12}$ теория не постулирует новых уровней; предельный
континуум содержит все допустимые структурные возможности.

% ================================================================
\subsubsection{Связь с M-пространствами}

$M_{12}$ определяет:
\begin{itemize}
    \item полную область допустимости для предельных структур,
    \item требования инвариантности для операторов,
    \item ограничения глобальной когерентности,
    \item границы замыкания ветвления,
    \item условия существования неподвижной точки,
    \item совместимость со всеми $M_0\dots M_{11}$.
\end{itemize}

Континуум существует на уровне $K_{12}$ тогда и только тогда, когда вся его
структура лежит внутри допустимой области $M_{12}$.
